\section{Legal Discussion}
\label{sec:legal}

\subsection{Scope of Claims}
\label{sec:legal:scope}

This paper demonstrates that dual-constrained bisection can reduce partisan
gaps in states with favourable Lorenz geometry (Wisconsin is the clearest
example).
It does \textbf{not} demonstrate:
\begin{enumerate}
  \item that courts can order this method to be used in redistricting;
  \item that any particular partisan outcome is constitutionally required;
  \item that the algorithm is partisan-neutral --- it explicitly incorporates
    a proportionality target and thus requires partisan vote data as an input;
  \item that states with $\sigma < 0.05$ (free proportionality) are presently
    gerrymandered; or
  \item that B.12 at any $\eta$ is guaranteed to improve proportionality ---
    Table~\ref{tab:tradeoff-empirical} shows that it worsens outcomes for North
    Carolina at all tested $\eta$ values.
\end{enumerate}
The appropriate use of these results is to characterise the \emph{cost in
partisan information} of achieving proportionality, and to distinguish states
where proportionality is a geographic impossibility (expensive proportionality
states) from states where it is feasible but unreached.

\subsection{Federal Prohibition (\S104(e))}

The hypothetical Districting Integrity Act \S104(e) prohibits partisan data
as an input to federal congressional redistricting.
B.12 is not a proposal to violate this prohibition.
Rather, it is a mathematical analysis of what proportionality would cost if
the prohibition were relaxed --- and a characterisation of when it costs nothing
(free proportionality states) versus when it costs significantly.

The \S104(e) prohibition is a policy choice, not a mathematical necessity.
The choice has consequences: in expensive-proportionality states, geography-only
algorithms produce maps where 50\,\% of voters elect 35\,\% of representatives.
B.12 measures those consequences.

\subsection{State Legislative Applications}

State legislatures are not constrained by \S104(e).
Several state constitutions require proportional representation or prohibit
maps that unduly favour a political party \citep{lwv2018pa, harper2022nc}.

For state legislative redistricting, B.12 provides a rigorous operationalisation
of proportionality: set $\mathtt{ubvec}[1]$ to the minimum value that achieves
proportional seat allocation, and verify the Lorenz feasibility condition is
satisfied.
If Lorenz feasibility fails (expensive-proportionality state), the court can
be informed that geographic constraints make proportionality unachievable
with single-member districts --- a fact about geography, not about intent.

\subsection{The Evidentiary Role}

B.12's most immediate legal contribution is evidentiary.
For states where $\sigma$ is low (free or cheap proportionality), a court
can assess whether a geography-only algorithm's non-proportional outcome is
``natural'' (unavoidable given geography) or artificially produced.
If $\sigma < 0.05$ and the algorithm still produces a $-10$\,pp gap,
something other than geography is causing the shortfall.

For states where Lorenz feasibility is violated (expensive-proportionality),
the court has a rigorous basis for finding that proportionality is geometrically
impossible with contiguous single-member districts, which may support a
different remedy (e.g., multi-member districts).

\subsection{\emph{Rucho v.\ Common Cause} (2019) and Algorithmic Neutrality}

\emph{Rucho v.\ Common Cause}, 588 U.S. 684 (2019), held that federal courts
cannot remedy partisan gerrymandering under the federal Constitution because
``partisan gerrymandering claims present political questions beyond the reach
of the federal courts.''
B.12's analysis does not challenge this holding, but it reframes what the
holding leaves open.

\emph{Rucho} was decided on justiciability grounds, not on the merits.
Chief Justice Roberts noted that ``[t]he courts have no license to reallocate
political power between the two major political parties.''
This reasoning applies directly to human redistricting decisions with
identifiable partisan intent.
B.12's algorithmic approach offers a different posture: an algorithm that
takes only census population and geographic adjacency as inputs (the
geography-only baseline) has \emph{no partisan intent} --- it cannot
gerrymander because it lacks the necessary mental state.
The \emph{Rucho} prohibition on judicial reallocations of political power
presupposes that political power has been deliberately allocated;
a purely algorithmic plan reallocates nothing.

In the post-\emph{Rucho} landscape, three legal developments matter for B.12:

\begin{enumerate}
  \item \textbf{State courts}: \emph{Rucho} explicitly invited state-court
    review under state constitutions. \emph{League of Women Voters v.\ Commonwealth
    of Pennsylvania} (Pa. 2018) and \emph{Harper v.\ Hall} (N.C. 2022) struck
    down maps as partisan gerrymanders under state constitutional anti-partisan
    provisions.
    B.12's $\sigma$ classification provides a rigorous quantitative test that
    state courts can apply: a map drawn by a geography-only algorithm with
    $\sigma = 0$ has no partisan signal and cannot be a partisan gerrymander.

  \item \textbf{The algorithmic neutrality standard}: B.12 defines a
    \emph{spectrum} from full algorithmic neutrality ($\sigma = 0$) to
    full partisan optimisation.
    The \emph{Rucho} Court observed that ``an unbiased redistricting system would
    produce a roughly proportional result.''
    B.12's proportionality paradox shows that this observation is geographically
    incorrect for competitive states: an unbiased system (B.9) does \emph{not}
    produce proportional results because of the Rodden effect.
    This factual finding is directly relevant to state court proportionality
    challenges: it establishes that non-proportional outcomes from neutral
    algorithms are a geographic inevitability, not evidence of partisan intent.

  \item \textbf{Procedural fairness as the justiciable claim}: B.12 supports
    a procedural-fairness argument that \emph{Rucho} did not foreclose.
    Even if \emph{Rucho} bars courts from mandating proportional outcomes,
    courts retain authority to mandate fair \emph{processes}.
    A requirement that redistricting use a publicly auditable, partisan-data-free
    algorithm (such as B.9 or a verified $\sigma = 0$ variant of B.12) is a
    procedural mandate, not an outcome mandate.
    State courts in Pennsylvania and North Carolina have accepted this framing.
\end{enumerate}

B.12's central legal contribution is thus not to challenge \emph{Rucho} but to
provide the mathematical infrastructure for the post-\emph{Rucho} state-court
and legislative landscape: a precise, auditable definition of partisan signal
strength that separates geographic inevitability ($\Delta$ caused by Rodden
effect) from partisan design ($\Delta$ caused by deliberate targeting).
